\section{Complete near-group branch in dimension twenty-four}
\label{sec:ng24-actions}
Consider a fusion category $\mathcal D$ of type $(1^8,4)$ and a categorical
$C_2$ action whose nontrivial element fixes precisely two invertible simples.
Then its equivariantization has based ring $R9\_1\_12$ or $R9\_1\_108$.
More precisely, the former occurs when the invertibles of $\mathcal D$ form
$D_8$, and the latter when they form $Q_8$. In particular this branch cannot
produce $R9\_1\_62$ or $R9\_1\_87$.

This is an exhaustive statement about complex fusion categories and actions.
Indeed dimension $24=2^3\cdot3$ implies solvability by
\cite[Theorem~1.6]{ENOSolvable}, hence weak group-theoreticity and
unitarizability by \cite[Theorem~5.20]{GHR}.
Reutter's Theorems~1--2 and Corollary~3~\cite{Reutter} allow tensor
autoequivalences and their coherence to be transported to unitary ones.
The near-group rules are
$g\rho=\rho=\rho g$ and $\rho^2=\sum_{g\in G}g+2\rho$.
Izumi's classification \cite[Theorems~1.1 and~6.1]{Izumi} shows
$G=D_8$ or $Q_8$, and supplies exactly three
unitary categories for each group. The explicit admissible tuples below
cover all six. Thus a unitary classification is being used only after
unitarizability has been justified.

\subsection{Normal form for all tensor autoequivalences}
We use Izumi, \emph{A Cuntz algebra approach to the classification of
near-group categories}, arXiv:1512.04288v2, Remark~3.6, Definition~4.8,
Theorem~4.9, and Lemma~6.13~\cite{Izumi}. Let $\pi$ be the unique irreducible
two-dimensional representation of $G$, on $K=\mathbb C^2$. Let $j$ be an
antiunitary commuting with $\pi(G)$ and satisfying $j^2=\epsilon$, with
$\epsilon=+1$ for $D_8$ and $\epsilon=-1$ for $Q_8$. For any third root
of unity $\zeta$, the complete categorical data are
\[
(K,V,U_K,j_1,j_2,\chi,l)
=(\mathbb C^2,\pi,\pi,j,\epsilon\zeta j,\chi,0),
\quad
\pi(g)\pi(h)=\chi_h(g)\pi(h)\pi(g).
\]
The generating fusion isometries are
$S_g\in\operatorname{Hom}(g,\rho^2)$ and
$T(v)\in\operatorname{Hom}(\rho,\rho^2)$, $v\in K$.
The left absorption $g\rho=\rho$ is strict in this presentation.
All morphisms are generated by these isometries and their adjoints
(Remark~3.6).

Every tensor autoequivalence is, up to tensor natural isomorphism,
represented by a group automorphism $\phi$ and a unitary matrix $W$ with
\begin{equation}\label{eq:ng-W}
 W\pi(g)W^*=\pi(\phi(g)),\qquad Wj=jW.
\end{equation}
The representative acts by
\begin{equation}\label{eq:ng-gamma}
 \gamma(S_g)=S_{\phi(g)},\qquad
 \gamma(T(v))=T(Wv).
\end{equation}
Here the normal-form assertion includes all tensorators, not just
permutations of simples. To see this, normalize the one-dimensional left
absorption maps $g\rho\to\rho$ under the functor. The tensor-functor
coherence on $(g,h,\rho)$ then normalizes the pointed multiplication
maps as well. More explicitly, if $\lambda_g$ is the tensorator on the
strict left absorption $g\rho=\rho$, the natural-isomorphism component
on $g$ makes $\lambda_g=1$, and the same coherence forces all pointed
tensorators to one. A further scalar choice on $\rho$ fixes $S_1$.
For a functor with tensorator $L$, its normalized trivalent images are
$L_{\rho,\rho}^{-1}F(S_g)$ and $L_{\rho,\rho}^{-1}F(T)$.
The following formulas refer to these normalized images.
Since $S_g$ is obtained by applying the left $g$ action to $S_1$,
its image is now $S_{\phi(g)}$. The action on $K$ is a unitary $W$.
Covariance of the left $G$ action and of the duality defined by $S_1$
gives exactly~\eqref{eq:ng-W}. Remark~3.6 then determines the functor on
all morphisms. This is also the equivalence-of-tuples normal form in
Definition~4.8 and the proof of Theorem~4.9.

Conversely~\eqref{eq:ng-W} preserves $\chi$ and both $j_1,j_2$, and
$l=0$ supplies no further constraint. The reconstruction formulas of
Section 4 therefore give the tensor autoequivalence~\eqref{eq:ng-gamma}.
Explicitly, the formula for $\rho(S_1)$ uses only
$\sum_g S_g$ and $\sum_iT_i j(T_i)$, both covariant under $\gamma$;
$\rho(S_g)=U(g)\rho(S_1)U(g)^*$ is also covariant. Formula (3.12) for
$\rho(T)$ is covariant because $W$ intertwines $\pi$ and commutes with
$j_1,j_2$. Thus this construction supplies genuine functors.

For each $\phi$, Schur's lemma gives $W$ uniquely up to a scalar. The
condition $Wj=jW$ reduces that ambiguity to $\pm W$. Their quotient is
the functor
\[
\delta(S_g)=S_g,\qquad\delta(T)=-T.
\]
It is tensor-isomorphic to the identity: take $\eta_g=1$ and
$\eta_\rho=-1$, extend tensorially, and check naturality on $S_g,T$.
For example $\eta_{\rho^2}\delta(T)=T\eta_\rho=-T$, while the
$S_g$ equation has both sides $S_g$. These generators suffice for all
naturality equations. Consequently there is no additional invisible
autoequivalence or free tensorator parameter at a fixed $\phi$.

Finally $\operatorname{Aut}_\otimes(\operatorname{Id}_{\mathcal D})=1$.
The relation $g\rho=\rho$ forces the scalar on every $g$ to be one,
and the occurrence of $\rho$ in $\rho^2$ forces the scalar on $\rho$
to be one. Thus a square isomorphism for an involutive tensor functor,
when it exists, is unique, and automatically satisfies the $C_2$ action
coherence. In particular there is no independent sign choice that could
alter the duality calculation below.

\subsection{The square of the action determines the large-object duality}
If $\phi^2=1$, then $W^2$ commutes with $\pi(G)$ and is scalar.
The antiunitary condition makes this scalar $+1$ or $-1$.
For $D_8$ choose the real matrices
\[
 \pi(r)=\begin{pmatrix}0&-1\\1&0\end{pmatrix},\quad
 \pi(s)=\begin{pmatrix}1&0\\0&-1\end{pmatrix},\quad j=\text{complex conjugation}.
\]
An automorphism with exactly two fixed group elements is
$r\mapsto r^{-1},s\mapsto r^k s$, $k=1$ or $3$.
Its intertwiner $W$ is a real reflection, so $\operatorname{tr}W=0$
and $W^2=1$. For $k=1$ one may take
$W=2^{-1/2}\left(\begin{smallmatrix}1&1\\1&-1\end{smallmatrix}\right)$.
The action is strict and its square coherence has $\eta_\rho=+1$.

For $Q_8$ use
\[
 \pi(r)=\begin{pmatrix}i&0\\0&-i\end{pmatrix},\quad
 \pi(s)=\begin{pmatrix}0&1\\-1&0\end{pmatrix},\quad
 j=\pi(s)\circ\text{complex conjugation}.
\]
The unitaries commuting with $j$ form $SU(2)$. Every nontrivial
involutive conjugation of $\pi(Q_8)$ lifts to a noncentral $W\in SU(2)$
with $W^2=-1$ and trace zero. In particular this holds for all six
involutions fixing precisely two group elements. Thus $\gamma^2=\delta$,
and the unique square coherence is $\eta_g=1$, $\eta_\rho=-1$.
It is invariant under $\gamma$, proving the action pentagon explicitly.

An equivariant structure on $\rho$ is a scalar $u$ with
$u^2=\eta_\rho$. Since $\gamma(S_1)=S_1$, the lift of $\rho$ is
self-dual exactly when $u^2=1$. Consequently the two lifts are both
self-dual for $D_8$, and form a dual pair for $Q_8$.
This conclusion is independent of the third-root parameter $\zeta$.

\subsection{Independent finite tensor replay}
The pointed part is the strict equivariantization of $\operatorname{Vec}_G$
by $\phi$. Its simple objects are the two characters over each fixed
$g\in\{1,r^2\}$ and one object for each of the three free orbits.
Products are obtained by counting ordered pairs $(g,h)$ with prescribed
product, and projecting the fixed output fibers onto the two $C_2$
characters. This computation allows noncommutative $G$ throughout.
Write $P_+,P_-$ for the two lifts of $\rho$. A fixed invertible
$(g,t)$ multiplies $P_u$ to $P_{tu}$, and each free-orbit simple multiplies
it to $P_++P_-$. For $P_tP_u$, each free invertible orbit occurs once,
and the fixed outputs are
$(1,\epsilon tu)$ and $(r^2,\epsilon tu)$.
Both $P_+$ and $P_-$ occur once because the action on
$\operatorname{Hom}(\rho,\rho^2)$ is $W$ with trace zero.
These rules determine every coefficient.

The standard-library verifier \nolinkurl{verify_ng_actions.py} enumerates
all $8$ automorphisms of $D_8$ and all $24$ of $Q_8$, selecting respectively
the $2$ and $6$ required involutions. For every involution it supplies
an exact matrix $W=A/\sqrt2$, with Gaussian-integer $A$, and checks
unitarity, the intertwining equations, antiunitary compatibility, trace
zero, and $W^2=\epsilon$. It then constructs the complete equivariant
fusion tensor and checks all associativity identities. Every $D_8$
result matches $R9\_1\_12$ and every $Q_8$ result matches
$R9\_1\_108$, coefficient by coefficient, with an explicit
model-to-input permutation. None matches $R9\_1\_62$ or $R9\_1\_87$.
The complete receipt is \nolinkurl{ng_actions.json}.

This argument resolves the entire $(1^8,4)$ de-equivariantization branch;
it does not by itself dispose of other possible de-equivariantization
profiles of a dimension-forty-eight input.
